\chapter{Flutter UI Widget Dictionary}
\label{ch:flutter-ui-widget-dictionary}

This appendix provides a comprehensive reference to the primary Flutter UI widgets and the design system implemented for the Kemora application. The frontend architecture relies on a highly modular and reusable component strategy, leveraging the ``Modern Heritage'' design aesthetic.

\section{Modern Heritage Design System}
\label{sec:modern-heritage-design-system}

The ``Modern Heritage'' design system is implemented in Flutter using Material 3's \texttt{ThemeData} and custom token classes (\texttt{AppColors}, \texttt{AppTypography}, and \texttt{AppShadows}). It combines rich, earthy tones like \texttt{primaryContainer} (\#c5a021, Primary Gold) and \texttt{secondary} (\#435b9f, Primary Blue) with modern, clean typography and subtle drop shadows for an elegant, premium feel.

\begin{itemize}
    \item \textbf{Color Scheme:} The \texttt{AppColors} class defines all color tokens, bridging semantic naming (e.g., \texttt{primary}, \texttt{surfaceContainerHighest}) with fixed brand colors.
    \item \textbf{Typography:} \texttt{AppTypography} extends Material 3's \texttt{TextTheme}, applying brand-specific font weights and letter spacing (tracking) to establish hierarchy.
    \item \textbf{Theme Injection:} \texttt{AppTheme.lightTheme} configures the global \texttt{ThemeData}, overriding default styles for \texttt{AppBar}, \texttt{ElevatedButton}, \texttt{Card}, and \texttt{InputDecoration} to ensure consistency without redundant inline styling.
\end{itemize}

\section{Core Reusable Widgets}
\label{sec:core-reusable-widgets}

The \texttt{lib/presentation/widgets} directory contains all globally shared UI components. These components encapsulate complex layout, state, and animation logic.

\subsection{EditorialPlaceCard}
The \texttt{EditorialPlaceCard} is a visually striking component used to display points of interest, accommodations, and tours. It features a large cover image, a soft gradient overlay for text legibility, and an integrated \texttt{GlassmorphismContainer} for the favorite action.

\begin{lstlisting}[language=Java, caption={EditorialPlaceCard Build Method Excerpt}]
  @override
  Widget build(BuildContext context) {
    return GestureDetector(
      onTap: onTap,
      child: AspectRatio(
        aspectRatio: aspectRatio > 1 ? aspectRatio : 1 / (1 / aspectRatio),
        child: Container(
          decoration: BoxDecoration(
            borderRadius: BorderRadius.circular(24),
            boxShadow: AppShadows.ambient,
          ),
          child: ClipRRect(
            borderRadius: BorderRadius.circular(24),
            child: Stack(
              fit: StackFit.expand,
              children: [
                // Image or Placeholder
                // ... inline image loading logic ...
                
                // Gradient Overlay (Soft Lapis-to-transparent)
                Container(
                  decoration: BoxDecoration(
                    gradient: LinearGradient(
                      begin: Alignment.bottomCenter,
                      end: Alignment.topCenter,
                      colors: [
                        AppColors.secondary.withOpacity(0.9),
                        AppColors.secondary.withOpacity(0.4),
                        Colors.transparent,
                      ],
                      stops: const [0.0, 0.4, 0.8],
                    ),
                  ),
                ),
                
                // Content Overlay ...
                // ... inline widget composition ...
              ],
            ),
          ),
        ),
      ),
    );
  }
\end{lstlisting}

\subsection{Itinerary Widgets}
The itinerary timeline is composed of several modular widgets defined in \texttt{itinerary\_widgets.dart}, including \texttt{DayHeader}, \texttt{ActivityTimelineTile}, and \texttt{PremiumActivityCard}.

The \texttt{ActivityTimelineTile} constructs the vertical timeline layout, seamlessly connecting adjacent activities using custom-drawn lines and circular nodes.

\begin{lstlisting}[language=Java, caption={ActivityTimelineTile Layout Structure}]
class ActivityTimelineTile extends StatelessWidget {
  final ItineraryItem activity;
  final bool isLast;
  const ActivityTimelineTile({
    super.key,
    required this.activity,
    this.isLast = false,
  });

  @override
  Widget build(BuildContext context) {
    return IntrinsicHeight(
      child: Row(
        crossAxisAlignment: CrossAxisAlignment.stretch,
        children: [
          // Timeline logic
          Column(
            children: [
              Container(
                width: 16,
                height: 16,
                decoration: BoxDecoration(
                  shape: BoxShape.circle,
                  color: AppColors.primaryContainer,
                  border: Border.all(color: Colors.white, width: 3),
                  boxShadow: [
                    BoxShadow(
                      color: AppColors.primaryContainer.withValues(alpha: 0.3),
                      blurRadius: 4,
                      offset: const Offset(0, 2),
                    ),
                  ],
                ),
              ),
              if (!isLast)
                Expanded(
                  child: Container(
                    width: 2,
                    color: AppColors.primaryContainer,
                    margin: const EdgeInsets.symmetric(vertical: 4),
                  ),
                ),
            ],
          ),
          const SizedBox(width: 16),
          // Content Card
          Expanded(
            child: Padding(
              padding: const EdgeInsets.only(bottom: 24),
              child: PremiumActivityCard(activity: activity),
            ),
          ),
        ],
      ),
    );
  }
}
\end{lstlisting}

For the full component dictionary, please refer to the GitHub repository.
